% =====================================================================
%  lang-el.tex — η ελληνική έκδοση.
%
%  Η ελληνική ΑΚΟΛΟΥΘΕΙ την αγγλική: το content/en/ είναι το πρωτότυπο,
%  το content/el/ η μετάφρασή του. Το tools/sync-check.py δείχνει ποια
%  ελληνικά κομμάτια έχουν μείνει πίσω. Δες docs/SYNC.md.
% =====================================================================
\setmainlanguage[numerals=arabic]{greek}
\setotherlanguage{english}

% Ελληνικές γραμματοσειρές: ίδιες με τις λατινικές, ώστε οι δύο εκδόσεις
% να μοιάζουν με το ίδιο έγγραφο.
\newfontfamily\greekfont[Ligatures=TeX]{\thesisserif}
\newfontfamily\greekfontsf[Ligatures=TeX]{\thesissans}
\newfontfamily\greekfonttt[Scale=0.90,Ligatures=TeX]{\thesismono}

% \en{...} : αγγλικός όρος μέσα σε ελληνικό κείμενο, με σωστό συλλαβισμό
\newcommand{\en}[1]{\textenglish{\emph{#1}}}
% \enup{...}: αγγλικό όνομα χωρίς πλάγια (Paxos, HotStuff, …)
\newcommand{\enup}[1]{\textenglish{#1}}
\newcommand{\gr}[1]{#1}

% Ελληνική σύμβαση αριθμών στα σχήματα: υποδιαστολή κόμμα.
\pgfkeys{/pgf/number format/.cd, use comma, 1000 sep={.}}

\newcommand{\dgblankleaftext}{(Κενό φύλλο)}

\newcommand{\thesisnames}{%
  \renewcommand{\chaptername}{ΚΕΦΑΛΑΙΟ}%
  \renewcommand{\figurename}{Σχήμα}%
  \renewcommand{\tablename}{Πίνακας}%
  \renewcommand{\bibname}{ΒΙΒΛΙΟΓΡΑΦΙΑ}%
  \renewcommand{\contentsname}{ΠΕΡΙΕΧΟΜΕΝΑ}%
  \renewcommand{\listfigurename}{ΚΑΤΑΛΟΓΟΣ ΣΧΗΜΑΤΩΝ}%
  \renewcommand{\listtablename}{ΚΑΤΑΛΟΓΟΣ ΠΙΝΑΚΩΝ}%
}
